% ============================================================================
% GALOSH — ベンチマーク表（日本語キャプション版）
% 必要: \usepackage{booktabs}
% 凡例: 太字 = blind/training-free 手法中で最良、\underline = 同2位。
%       Blind2Unblind は学習済みネットワークで上限参照としてのみ掲載。
% ============================================================================

% ---------------------------------------------------------------- SIDD Medium
\begin{table}[t]
\centering
\caption{\textbf{SIDD Medium}（80 枚のフル解像度実ノイズ Bayer 画像）における
raw 領域ブラインド・デノイジング。全手法が \emph{blind}（クリーン/ノイズ対なし、
画像ごとのノイズオラクルなし）かつ \emph{training-free}（Blind2Unblind を除く）。
太字 / \underline{下線} = blind・training-free 手法中の最良 / 2位。
$\uparrow$: 高いほど良い、$\downarrow$: 低いほど良い。}
\label{tab:sidd_medium}
\setlength{\tabcolsep}{2pt}
\footnotesize
\resizebox{\columnwidth}{!}{%
\begin{tabular}{l ccccc rr}
\toprule
 & & & & & & \multicolumn{2}{c}{時間\,(s)}\\
\cmidrule(lr){7-8}
手法 & PSNR$\uparrow$ & SSIM$\uparrow$ & LPIPS$\downarrow$ & DISTS$\downarrow$ & NIQE$\downarrow$ & CPU & GPU\\
\midrule
\multicolumn{8}{l}{\textit{GALOSH（提案）--- blind・training-free}}\\
\quad GALOSH FP32          & \textbf{48.13} & \textbf{0.9883} & \textbf{0.2030} & \textbf{0.1690} & 8.62 & 12.6 & 0.76\\
\quad GALOSH INT16\,$^{\dagger}$ & \underline{47.42} & \underline{0.9881} & \underline{0.2030} & \underline{0.1694} & \underline{8.62} & 103.6$^{\dagger}$ & 2.7\\
% GALOSH 数値 = canonical GPU パイプライン（jinc アンチリンギング）、SIDD-Medium 全80枚:
% FP32 = o32 GPU、INT16 = i16 GPU（lf=5 cf=9。r32 CPU 参照に対し near-lossless
% ~58-65 dB end-to-end（フルフレーム）= INT16 ストレージ量子化そのもの。
% INT32 ストレージなら GPU は r32 と bit-exact、2026-07-04 検証済み）。
% galosh_fp32 + vst_galosh は 2026-07-02 に GPU exe で再検証（_metrics_revst_gpu.json、失敗0）。
% GPU 時間: fp32/vst = 同一セッションの専用非競合計測、int16 = i16 campaign 平均。
\addlinespace[2pt]
\multicolumn{8}{l}{\textit{アブレーション --- 共通 VST 前段（ネイティブ推定なし）}}\\
\quad VST\,+\,GALOSH core   & 44.35 & 0.9765 & 0.2132 & 0.1761 & \textbf{8.57} & 15.3 & 0.74\\
\quad VST\,+\,BM3D-CFA      & 45.96 & 0.9851 & 0.2932 & 0.2778 & 11.36 & 40.5 & ---\\
\quad VST\,+\,NLM-CFA       & 40.36 & 0.9679 & 0.4048 & 0.4194 & 12.98 & --- & 1.6\\
\addlinespace[2pt]
\multicolumn{8}{l}{\textit{古典的手法 --- blind、ネイティブ MAD}}\\
\quad BM3D-CFA             & 46.74 & 0.9862 & 0.2725 & 0.2613 & 10.55 & 40.4 & ---\\
\quad NLM-CFA              & 42.32 & 0.9729 & 0.3791 & 0.3901 & 12.05 & 27.1 & 1.2\\
\midrule
\multicolumn{8}{l}{\textit{学習済み DL（上限参照、blind ではない）}}\\
\quad Blind2Unblind        & 49.07 & 0.9924 & 0.1169 & 0.1281 & 9.43 & --- & 5.6\\
\bottomrule
\end{tabular}}

\vspace{2pt}
{\footnotesize 画像あたりのフルサイズ時間をプラットフォーム別の列に分離
（同じ列同士のみ比較可能）。CPU 列はすべてシングルスレッド実装。「---」= その
プラットフォームでは未計測（BM3D に GPU 移植はなく、Blind2Unblind は CUDA
必須）。NLM-CFA：GPU = CUDA 実装、CPU = 同一 patch/search パラメータの skimage
参照実装（フルフレーム 12 枚の平均；RawNIND の CPU 値は 100 シーン標本の中央値）。GALOSH は両列に値を持つ
--- クロスプラットフォームの基準点である。$^{\dagger}$INT16 の品質はストリーミング
INT16 GPU パイプラインで直接測定した値。CPU 時間は正確性優先の INT32 参照実装
（INT16 出力に対し near-lossless、速度未最適化）、GPU 時間はストリーミング
i16 実装（\ref{sec:map}節参照）。}
\end{table}

% ---------------------------------------------------------------- RawNIND
\begin{table}[t]
\centering
\caption{\textbf{RawNIND}（1493 枚のカメラ横断 $512{\times}512$ 実ノイズクロップ）に
おける raw 領域ブラインド・デノイジング。プロトコルと強調は
表~\ref{tab:sidd_medium} と同じ。}
\label{tab:rawnind}
\setlength{\tabcolsep}{2pt}
\footnotesize
\resizebox{\columnwidth}{!}{%
\begin{tabular}{l ccccc rr}
\toprule
 & & & & & & \multicolumn{2}{c}{時間\,(s)}\\
\cmidrule(lr){7-8}
手法 & PSNR$\uparrow$ & SSIM$\uparrow$ & LPIPS$\downarrow$ & DISTS$\downarrow$ & NIQE$\downarrow$ & CPU & GPU\\
\midrule
\multicolumn{8}{l}{\textit{GALOSH（提案）--- blind・training-free}}\\
\quad GALOSH FP32          & \underline{30.49} & \underline{0.7908} & \textbf{0.2395} & \textbf{0.2127} & \underline{8.64} & 0.54 & 0.45\\
\quad GALOSH INT16\,$^{\dagger}$ & 30.17 & 0.7860 & \underline{0.2453} & \underline{0.2161} & \textbf{8.63} & 2.32$^{\dagger}$ & 0.35\\
\addlinespace[2pt]
\multicolumn{8}{l}{\textit{アブレーション --- 共通 VST 前段（ネイティブ推定なし）}}\\
\quad VST\,+\,GALOSH core   & \textbf{30.63} & \textbf{0.8163} & 0.2481 & 0.2198 & 8.65 & 0.49 & 0.60\\
% vst_galosh は 2026-07-02 に GPU exe（外部 α/σ² 対応後）で再実行：現行 anti-ringing
% パイプライン、1493/1493、err 0（ベンチ JSON から機械照合。
% benchmark/scripts/verify_table_numbers.py 参照）。
% GPU 時間 fp32/vst = 専用非競合計測（_revst_gpu_time_rawnind.json、150 シーン標本）。
\quad VST\,+\,BM3D-CFA      & 30.07 & 0.7792 & 0.4320 & 0.3563 & 13.04 & 2.29 & ---\\
\quad VST\,+\,NLM-CFA       & 29.03 & 0.7649 & 0.5452 & 0.4636 & 14.57 & --- & 0.68\\
\addlinespace[2pt]
\multicolumn{8}{l}{\textit{古典的手法 --- blind、ネイティブ MAD}}\\
\quad BM3D-CFA             & 30.28 & 0.7848 & 0.4268 & 0.3461 & 12.15 & 2.30 & ---\\
\quad NLM-CFA              & 29.54 & 0.7735 & 0.5190 & 0.4364 & 13.80 & 0.58 & 0.76\\
\midrule
\multicolumn{8}{l}{\textit{学習済み DL（上限参照、blind ではない）}}\\
\quad Blind2Unblind        & 30.68 & 0.7932 & 0.2215 & 0.2018 & 8.93 & --- & ---\\
\bottomrule
\end{tabular}}

\vspace{2pt}
{\footnotesize プラットフォーム分離は表~\ref{tab:sidd_medium} と同じ。$512^2$
では GPU の起動/同期オーバーヘッドが GALOSH で支配的になるため、CPU ビルドが
ここでは競争的である。$^{\dagger}$表~\ref{tab:sidd_medium} 参照。Blind2Unblind
の計時は事前キャッシュ読み込みのため省略。}
\end{table}

% ------------------------------------------- Speed（統制参考）
\begin{table}[t]
\centering
\caption{\textbf{統制解像度の速度参考値}：1080p（1920$\times$1080 Bayer）、
CPU（1 コア）と GPU（NVIDIA RTX 4070\,Ti）。速度の\emph{主比較}は
表~\ref{tab:sidd_medium}--\ref{tab:rawnind} のフルサイズ画像あたり CPU/GPU 列で
ある。ウォーム実測。GALOSH GPU はカーネルパイプライン（context+build
$\sim$90\,ms は初回のみで償却）、他は end-to-end。}
\label{tab:speed}
\setlength{\tabcolsep}{8pt}
\resizebox{\columnwidth}{!}{%
\begin{tabular}{l rr}
\toprule
手法 & CPU (1 コア) & GPU\\
\midrule
\multicolumn{3}{l}{\textit{GALOSH（提案）--- blind・training-free}}\\
\quad GALOSH FP32 (o / o32)        & 1.38 s & \textbf{18.2 ms}\\
\midrule
\multicolumn{3}{l}{\textit{古典的手法（blind）}}\\
\quad BM3D-CFA                     & 6.33 s & ---\\
\quad NLM-CFA                      & 3.21 s & 0.42 s\\
\midrule
\multicolumn{3}{l}{\textit{学習済み DL}}\\
\quad Blind2Unblind                & ---$^{\dagger}$ & 0.74 s\\
\bottomrule
\end{tabular}}

\vspace{2pt}
{\footnotesize GALOSH FP32（GPU）は 4K で 46\,ms、5K で 72\,ms までスケールし
--- 同一 GPU 上で NLM-CFA の約 23$\times$、Blind2Unblind の約 40$\times$ ---
CPU でも古典的ベースラインの $2$--$5\times$ 高速である（1.38\,s 対
3.21/6.33\,s）：両プラットフォームで最速である。$^{\dagger}$Blind2Unblind の
参照実装は CUDA 必須で CPU では動作しない。INT16 固定小数点版
（\ref{sec:map}節）は CPU/GPU スループットではなく固定小数点ストリーミング
ハードウェアへの写像を狙うため、速度エントリには含めない。}
\end{table}
